\section*{Assessments and Final Grade}
In place of the familiar three-midterm-exams-and-a-final assessment structure, 
this course uses a Mastery-Based Learning (MBL) approach. 
MBL is a \href{https://drive.google.com/drive/folders/1-yut_EoMABRQYGjEdYB0vnV7DeqUL59f?usp=sharing}{research-backed} way 
to encourage learning from failure and 
help you track and prioritize your learning progress. 

\subsection*{List of Assessed Skills}
The MBL assessment approach divides course learning targets into individual skills. 
(Each MBL skill is roughly equivalent to one ``exam problem type'' in a traditional structure.) 
Figure~\ref{fig:skillsList} shows how these skills are divided into three categories: 
Must-Have (MH), Important (IMP), and Advanced (ADV). 
Each category includes both theory/paper skills and coding skills. 
This list does not include the \href{https://github.com/RHIT-CSSE/csse220/blob/master/Docs/debugging_skills_check.md}{Debugging Skill Check} (DSC), 
which will be assessed during drop-in hours or by appointment. 

\begin{figure}[!ht]
\centering
\includegraphics[width=0.85\textwidth]{skills-list}
\caption{MBL Skills List for CSSE 220.}%
\label{fig:skillsList}
\end{figure}

\subsection*{MBL Assessment Grading and Timing}
Throughout the course, 
you will have multiple attempts to demonstrate each skill. 
Each attempt will be graded as ``Met'' or ``Not Yet''. 
Once you earn a ``Met'' for a skill, 
you do not need to attempt that skill again 
(though you \emph{will} use that skill 
in future assignments, 
especially the final project). 

The theory/paper skills will be assessed 
during the second 50-minute class period, 
as indicated in the course calendar on Moodle. 
The coding skills will be assessed 
during the entire 110-minute class session, 
as indicated in the course calendar on Moodle. 
On each skill assessment day, 
you will select which of the available skills to attempt, 
which may include skills for which you received a ``Not Yet'' previously, 
by picking up the instructions sheets for those skills. 
See Figures~\ref{fig:codingSkillsSchedule} and~\ref{fig:paperSkillsSchedule} 
for the schedule of when each skill will be assessed for the first time. 

% This is where fig:codingSkillsSchedule and fig:paperSkillsSchedule should go

The final exam slot will be used as a last opportunity 
to attempt skills you have not yet met. 
Only one new skill, ADV03, will be introduced during the final exam slot. 
There is no limit on the number of skills 
you may attempt during the final exam slot, 
but it is in your best interest to meet 
as many skills as possible \emph{before} finals week. 

\subsection*{Final Letter Grades}
Table~\ref{tab:letterGradeReqs} shows 
the minimum requirements for each letter grade. 
(Failure to meet the minimum requirements for a D results in an F.) 
\textbf{We reserve the right to modify these requirements in a manner 
that makes some or all letter grades \emph{easier} to attain.} 
We will not make any requirements harder. 
\begingroup
\renewcommand{\arraystretch}{0.85}% reduce row height
\setlength{\tabcolsep}{6pt}% slightly tighter column padding
\begin{table}[!ht]
\caption{Requirements to earn each possible letter grade.}%
\label{tab:letterGradeReqs}
\begin{tabular}{lp{12cm}}
\toprule
\textbf{Final Grade} & \textbf{Minimum Requirements} \\
\midrule
D   &   \begin{itemize}[nolistsep,noitemsep]
        \item meet all but two MH skills
        \item pass Debugging Skill Check
        \item complete final project at ``satisfactory'' level or better
        \item complete $\ge 60\%$ of homework assignments
        \end{itemize}\\
\midrule
D+  &   Meet all requirements for D, but also %: 
        %\begin{itemize}[nolistsep,noitemsep]
        %\item 
        meet two more skills %(any level)
        %\end{itemize} 
        \\
\midrule
C   &   \begin{itemize}[nolistsep,noitemsep]
        \item meet all MH skills %\update{except one$^*$}
        \item pass Debugging Skill Check
        \item complete final project at ``satisfactory'' level or better
        \item complete $\ge 70\%$ of homework assignments
        \end{itemize} \\
\midrule
C+  &   Meet all requirements for C, but also %: 
        %\begin{itemize}[nolistsep,noitemsep]
        %\item 
        meet 1 IMP + 
        %\item meet one more IMP/ADV
        1 more IMP/ADV
        %\end{itemize}
        \\
\midrule
B   &   \begin{itemize}[nolistsep,noitemsep]
        \item meet all MH skills
        \item meet two IMP skills
        \item pass Debugging Skill Check
        \item complete final project at ``good'' level or better
        \item complete $\ge 75\%$ of homework assignments
        \end{itemize}\\
\midrule
B+  &   Meet all requirements for B, but also %:
        %\begin{itemize}[nolistsep,noitemsep]
        %\item 
        meet two more IMP/ADV skills
        %\end{itemize}
        \\
\midrule
A   &   \begin{itemize}[nolistsep,noitemsep]
        \item meet all MH and IMP skills
        \item meet one ADV skill
        \item pass Debugging Skill Check
        \item complete final project at ``good'' level or better
        \item complete $\ge 85\%$ of homework assignments
        \end{itemize}\\
\bottomrule
\end{tabular}
\end{table}
\endgroup

% \update{$^*$To be eligible for this exception, you must (1) have Met at least two IMP/ADV skills; (2) meet with me during Week 10 or finals week to review your Not Yet MH skills; \emph{and} (3) make an honest attempt on all Not Yet MH skills during the final exam slot.} 

\vspace{1em}

% Ugly hack to put these here, but it works for now
\begin{figure}[!ht]
\centering
\includegraphics[height=4cm]{coding-schedule}
\caption{Schedule of Coding Skill Assessments}%
\label{fig:codingSkillsSchedule}
\end{figure}

\begin{figure}[!ht]
\centering
\includegraphics[width=\textwidth]{paper-schedule}
\caption{Schedule of Theory/Paper Skill Assessments}%
\label{fig:paperSkillsSchedule}
\end{figure}
